\documentclass{article}

% packages
\usepackage{amsmath, amsthm, thmtools, amsfonts, amssymb, luacode, catchfile, tikzducks, hyperref, ifthen}
\ifcsname c@kobocompile\endcsname
	\usepackage[a5paper, total={1072pt, 1448pt}, margin=10pt, includeheadfoot]{geometry} % set page margins
\else
	\usepackage[a4paper, margin=50pt, includeheadfoot]{geometry}
\fi
\usepackage[shortlabels]{enumitem}
\usepackage[skip=3pt, indent=0pt]{parskip}

% language
\usepackage[bidi=basic, layout=tabular, provide=*]{babel}
\ifcsname c@english\endcsname
	\babelprovide[main, import]{english}
\else
	\babelprovide[main, import]{hebrew}
	\babelprovide{rl}
\fi
%\babelfont{rm}{Libertinus Serif}
\babelfont{rm}[Renderer=Harfbuzz]{Libertinus Serif}
\babelfont{sf}{Libertinus Sans}
\babelfont{tt}{Libertinus Mono}

% style
\AddToHook{cmd/section/before}{\clearpage}	% Add line break before section
\linespread{1.3}
\setcounter{secnumdepth}{0}		% Remove default number tags from sections, this won't do well with theorems
\AtBeginDocument{\setlength{\belowdisplayskip}{3pt}}
\AtBeginDocument{\setlength{\abovedisplayskip}{3pt}}
\graphicspath{ {../images/} }

% operators
\DeclareMathOperator\cis{cis}
\DeclareMathOperator\Sp{Sp}
\DeclareMathOperator\tr{tr}
\DeclareMathOperator\im{Im}
\DeclareMathOperator\re{Re}
\DeclareMathOperator\diag{diag}
\DeclareMathOperator*\lowlim{\underline{lim}}
\DeclareMathOperator*\uplim{\overline{lim}}
\DeclareMathOperator\rng{rng}
\DeclareMathOperator\Sym{Sym}
\DeclareMathOperator\Arg{Arg}
\DeclareMathOperator\Log{Log}
\DeclareMathOperator\dom{dom}
\DeclareMathOperator\supp{Supp}
\DeclareMathOperator\var{Var}
\DeclareMathOperator\cov{Cov}

% commands
%\renewcommand\qedsymbol{\textbf{מש''ל}}
%\renewcommand\qedsymbol{\fbox{\emoji{lizard}}}
\newcommand{\Aa}[0]{\mathcal{A}}
\newcommand{\Bb}[0]{\mathcal{B}}
\newcommand{\CC}[0]{\mathbb{C}}
\newcommand{\Cc}[0]{\mathcal{C}}
\newcommand{\EE}[0]{\mathbb{E}}
\newcommand{\FF}[0]{\mathbb{F}}
\newcommand{\Ff}[0]{\mathcal{F}}
\newcommand{\Ii}[0]{\mathcal{I}}
\newcommand{\Gg}[0]{\mathcal{G}}
\newcommand{\Ll}[0]{\mathcal{L}}
\newcommand{\Mm}[0]{\mathcal{M}}
\newcommand{\NN}[0]{\mathbb{N}}
\newcommand{\Nn}[0]{\mathcal{N}}
\newcommand{\PP}[0]{\mathbb{P}}
\newcommand{\Pp}[0]{\mathcal{P}}
\newcommand{\QQ}[0]{\mathbb{Q}}
\newcommand{\RR}[0]{\mathbb{R}}
\newcommand{\Rr}[0]{\mathcal{R}}
\newcommand{\Ss}[0]{\mathcal{S}}
\newcommand{\TT}[0]{\mathbb{T}}
\newcommand{\Uu}[0]{\mathcal{U}}
\newcommand{\Vv}[0]{\mathcal{V}}
\newcommand{\Ww}[0]{\mathcal{W}}
\newcommand{\ZZ}[0]{\mathbb{Z}}
\newcommand{\acts}[0]{\circlearrowright}
\newcommand{\explain}[2] {
	\begin{flalign*}
		 && \text{#2} && \text{#1}
	\end{flalign*}
}
\newcommand{\maketitleprint}[0]{ \begin{center}
	%\begin{tikzpicture}[scale=3]
	%	\duck[graduate=gray!20!black, tassel=red!70!black]
	%\end{tikzpicture}	
	\includegraphics[width=6cm]{cover}
\end{center}
}

% theorem commands
\newtheoremstyle{c_remark}
	{}	% Space above
	{}	% Space below
	{}% Body font
	{}	% Indent amount
	{\bfseries}	% Theorem head font
	{}	% Punctuation after theorem head
	{.5em}	% Space after theorem head
	{\thmname{#1}\thmnumber{ #2}\thmnote{ \normalfont{\text{(#3)}}}}	% head content
\newtheoremstyle{c_definition}
	{3pt}	% Space above
	{3pt}	% Space below
	{}% Body font
	{}	% Indent amount
	{\bfseries}	% Theorem head font
	{}	% Punctuation after theorem head
	{.5em}	% Space after theorem head
	{\thmname{#1}\thmnumber{ #2}\thmnote{ \normalfont{\text{(#3)}}}}	% head content
\newtheoremstyle{c_plain}
	{3pt}	% Space above
	{3pt}	% Space below
	{\itshape}% Body font
	{}	% Indent amount
	{\bfseries}	% Theorem head font
	{}	% Punctuation after theorem head
	{.5em}	% Space after theorem head
	{\thmname{#1}\thmnumber{ #2}\thmnote{ \text{(#3)}}}	% head content

\ifcsname c@english\endcsname
	\theoremstyle{plain}
	\newtheorem{theorem}{Theorem}[section]
	\newtheorem{lemma}[theorem]{Lemma}
	\newtheorem{proposition}[theorem]{Proposition}
	\newtheorem*{proposition*}{Proposition}
	%\newtheorem{corollary}[theorem]{אין חלופה עברית}

	\theoremstyle{definition}
	\newtheorem{definition}[theorem]{Definition}
	\newtheorem*{definition*}{Definition}
	\newtheorem{example}{Example}[section]
	\newtheorem{exercise}{Exercise}[section]

	\theoremstyle{remark}
	\newtheorem*{remark}{Remark}
	\newtheorem*{solution}{Solution}
	\newtheorem{conclusion}[theorem]{Conclusion}
	\newtheorem{notation}[theorem]{Notation}
\else
	\theoremstyle{c_plain}
	\newtheorem{theorem}{משפט}[section]
	\newtheorem{lemma}[theorem]{למה}
	\newtheorem{proposition}[theorem]{טענה}
	\newtheorem*{proposition*}{טענה}
	%\newtheorem{corollary}[theorem]{אין חלופה עברית}

	\theoremstyle{c_definition}
	\newtheorem{definition}[theorem]{הגדרה}
	\newtheorem*{definition*}{הגדרה}
	\newtheorem{example}{דוגמה}[section]
	\newtheorem{exercise}{תרגיל}[section]

	\theoremstyle{c_remark}
	\newtheorem*{remark}{הערה}
	\newtheorem*{solution}{פתרון}
	\newtheorem{conclusion}[theorem]{מסקנה}
	\newtheorem{notation}[theorem]{סימון}
\fi

% Questions related commands
\newcounter{question}
\setcounter{question}{1}
\newcounter{sub_question}
\setcounter{sub_question}{1}

\ifcsname c@english\endcsname
	\newcommand{\question}[1][0]{
		\ifthenelse{#1 = 0}{}{\setcounter{question}{#1}}
		\section{Question \arabic{question}}
		\addtocounter{question}{1}
		\setcounter{sub_question}{1}
	}

	\newcommand{\subquestion}[1][0]{
		\ifthenelse{#1 = 0}{}{\setcounter{sub_question}{#1}}
		\subsection{Part \alph{sub_question}}
		\addtocounter{sub_question}{1}
	}
\else
	\newcommand{\question}[1][0]{
		\ifthenelse{#1 = 0}{}{\setcounter{question}{#1}}
		\section{שאלה \arabic{question}}
		\addtocounter{question}{1}
		\setcounter{sub_question}{1}
	}

	\newcommand{\subquestion}[1][0]{
		\ifthenelse{#1 = 0}{}{\setcounter{sub_question}{#1}}
		\subsection{סעיף \localecounter{letters.gershayim}{sub_question}}
		\addtocounter{sub_question}{1}
	}
\fi

% import lua and start of document
\directlua{common = require ('../common')}

\GetEnv{AUTHOR}

% headers
\author{\AUTHOR}
\date\today

\title{פתרון מטלה 03 --- תורת הקבוצות האקסיומטית, 80650}

\DeclareMathOperator{\trcl}{trcl}
\DeclareMathOperator{\rank}{rank}

\begin{document}
\maketitle
\maketitleprint{}

\Question{}
נניח $\textbf{ZF} + \textbf{AC}$.
יהי מונה רגולרי אינסופי $\kappa$. נגדיר $H(\kappa) = \{ x \mid |\trcl(x)| < \kappa \}$  כאשר $\trcl(x)$ הסגור הטרנזיטיבי של $x$.

תהי קבוצה טרנזיטיבית $x$.
נוכיח ש־$z = \{\rank y \mid y \in x \}$ קבוצה טרנזיטיבית.
\begin{proof}
	נגדיר $c : z \to x$ פונקציית בחירה כך ש־$c(a) \in x$ וכמובן $\rank c(a) \in z$. \\*
	נוכיח שאם $t = \rank x$, אז לא קיים $l < t$ כך ש־$\forall y \in x, \rank y \ne l$. \\*
	נבחן את $l + 1 = t$ עבור $t$ עוקב. מההגדרה השקולה ל־$\rank$ מהמטלה הקודמת נובע ישירות כי  קיים $\rank y = l, y \in x$. אילו $t$ גבולי, אז מאותה סיבה לכל $l < t$ קיים $y$ כזה אחרת $l + 1$ לא בקבוצה וההגדרה הגבולית לא חלה. \\*
	באינדוקציה מהמבנה הזה ועל־ידי שימוש ב־$c$ נוכל לקבל שאכן לא קיים $y$ כזה. \\*
	לבסוף, אם $b \in a \in z$, אז בהכרח $\exists y \in x, \rank y = b$ ולכן גם $c(b) \in x \implies b \in z$.
\end{proof}

\Question{}
נוכיח ש־$H(\kappa)$ קבוצה.
\begin{proof}
	יהי $x \in H(\kappa)$, ונגדיר $y = \trcl(x)$, אז $|y| < \kappa$. \\*
	נגדיר $\alpha = \rank y$, ונניח $\alpha > \kappa^+$ אז נקבל ש־$|y| > \kappa$ וזו סתירה להנחה, לכן $\rank y < \kappa^+$. \\*
	בהתאם $y \in V_{\kappa^+}$ לכל $y$ ולכן $H(\kappa) \subseteq V_{\kappa^+}$ ובפרט קבוצה.
\end{proof}

\Question{}
נוכיח ש־$\kappa \subseteq H(\kappa)$ וש־$H(\kappa)$ טרנזיטיבי.
\begin{proof}
	יהי $\alpha \in \kappa$, אז $\alpha = \trcl(\alpha)$ ואנו יודעים ש־$\kappa$ מונה ולכן מהגדרה $|\alpha| < \kappa$, שאם לא כן המונה בו $\alpha$ היה מוכל היה גדול מ־$\kappa$, וכמובן $\kappa \notin \kappa$.
	נסיק אם כך ש־$\alpha \in H(\kappa)$ לכל $\alpha \in \kappa$ ולכן $\kappa \subseteq H(\kappa)$. \\*
	נראה ש־$H(\kappa)$ טרנזיטיבי.
	יהי $x \in H(\kappa)$, אז $x$ טרנזיטיבי, ונבחן את $y = \trcl(x)$. כמובן $\trcl(y) = y$ ולכן $y \in H(\kappa)$ ומצאנו ש־$H(\kappa)$ טרנזיטיבית עבור הקבוצות הטרנזיטיביות בה.
	יהי $z \in x$, אז גם $z \in y$ ולכן $z \subseteq y$ ובהתאם $y$ סגור טרנזיטיבי של $z$ ו־$z \in H(\kappa)$ ולכן זו קבוצה טרנזיטיבית.

	מתרגיל 1 והטרנזיטיביות נוכל להסיק ש־$\langle H(\kappa), \in \rangle$ מקיים את אקסיומת ההיקפיות והיסוד והקבוצה הריקה. \\*
	אילו נניח ש־$\omega < \kappa$ אז גם $\omega \in H(\kappa)$ ומתרגיל אחד מתקיימת אקסיומת האינסוף.
\end{proof}

\Question{}
נניח ש־$x \in H(\kappa)$ ו־$f : x \to H(\kappa)$.
נראה ש־$f, \rng f \in H(\kappa)$.
\begin{proof}
	נתון ש־$x \in H(\kappa) \iff |\trcl(x)| < \kappa$ אבל $|x| \le |\trcl(x)|$ ולכן $|x| < \kappa$, אז תהי $g : x \to \alpha$ עבור $\alpha \in \kappa$, חד־חד ערכית. \\*
	נגדיר פונקציית החלפה $c : \rng f \to x$, היא חד־חד ערכית מהגדרה, ולכן גם $c \circ g$ חד־חד ערכית ונסיק $|\rng f| = \alpha$. \\*
	אנו יודעים כי איחוד של קבוצות טרנזיטיביות הוא טרנזיטיבי ולכן $u = \bigcup_{y \in \rng f} \trcl y$ קבוצה טרנזיטיבית. \\*
	עוד נבחין כי $|u| < |\alpha \times \kappa| \le |\kappa \times \kappa| = \kappa$ ולכן $\rng f \in H(\kappa)$. \\*
	לבסוף $|x \times \rng f| < |\kappa \times \kappa| < \kappa$ ולכן גם $f \in H(\kappa)$.

	תחת המודל, נסיק שלכל $x \in H(k)$ ומחלקת פונקציה $F$, גם $\{F(y) \mid y \in x\}$ קבוצה, שכן $F \upharpoonright x : x \to H(\kappa)$. \\*
	יהי $x \in H(\kappa)$, אז מעיקרון הסדר הטוב קיים $\alpha \in Ord$ כך שקיימת $f : x \to \alpha$ הפיכה, ומבחירה נוכל להסיק שגם $\alpha < \kappa$.
\end{proof}

\Question{}
נוכיח ש־$H(\kappa)$ מספקת את אקסיומת האיחוד.
\begin{proof}
	למעשה כבר ראינו את הטענה הזו בסעיף הקודם.
	איחוד של קבוצות טרנזיטיביות הוא טרנזיטיבי ולכן אם $x \in H(\kappa)$ אז $\bigcup x \subseteq \bigcup \trcl(x)$ ולכן $\bigcup x \in H(\kappa)$. \\*
	נקבל אם כך שאקסיומת האיחוד חלה עבור $\langle H(\kappa), \in \rangle$ ולכן זהו מודל טרנזיטיבי של $\textbf{ZFC} - \textbf{Power set}$.
\end{proof}

\Question{}
נראה ש־$\textbf{ZFC} - \textbf{Power Set}$ לא מוכיח את הקיום של קבוצה שאיננה בת־מניה.
\begin{proof}
	נגדיר $\kappa = \omega$ ולכן $\langle H(\omega), \in \rangle$ מודל של $\textbf{ZFC} - \textbf{Power Set}$, אז זהו מודל כך שאין בו קבוצות שאינן־בנות מניה.

	אם נוסיף למודל את האקסיומה שאין קבוצה שלא בת־מניה נקבל שהיא עקבית.
\end{proof}

\end{document}
